\documentclass[12pt]{article}
\usepackage{amsmath,amssymb}
\usepackage[margin=1in]{geometry}
\usepackage{natbib}
\usepackage[skip=0.5\baselineskip,indent=0pt]{parskip}
\usepackage[T1]{fontenc}
\usepackage{mathpazo}
\usepackage{titlesec}
\renewcommand{\thesection}{\Roman{section}}
\titleformat{\section}{\normalfont\Large\bfseries}{\thesection}{0.5em}{}
\titleformat{\subsection}{\normalfont\large\bfseries}{\thesubsection}{0.5em}{}
\usepackage{enumitem}
\setlist{leftmargin=*,labelindent=0pt,labelsep=0.325em,align=left}
\setlist[enumerate]{label=\roman*.,ref=\roman*}
\renewenvironment{itemize}{\begin{enumerate}}{\end{enumerate}}
\setlist[description]{style=unboxed,leftmargin=0pt}
\titlespacing*{\section}{0pt}{\baselineskip}{\baselineskip}
\titlespacing*{\subsection}{0pt}{\baselineskip}{\baselineskip}

\title{\makebox[\textwidth][l]{\parbox{\dimexpr\textwidth+0.5in\relax}{\raggedright Review of Gregory Clark: \emph{To Have and Have Not:}\\\emph{The Determinants of Social Status in England, 1600--2026}}}}
\author{David Cesarini}
\date{}

\begin{document}
\maketitle

% =====================================================================
\section{Introduction}

The central question of this review is whether Clark's conclusions follow
from his evidence. Eight concerns organize the critique:
\begin{enumerate}
\item \textbf{Good fit does not establish a correct model.}
Even satisfactory fit need not identify the causal mechanism or yield
unbiased genetic estimates---the Goldberger point.

\item \textbf{Environmental transmission is excluded by assumption.}
The model makes shared genes the only transmitted source of resemblance
among relatives. Spousal resemblance induced by assortment requires a
qualification to that statement.

\item \textbf{Independent evidence challenges the assumptions.}
Assess evidence from twins, adoption, extended families, and molecular
genetics, specifying which restriction each finding challenges and
evaluating Clark's response.

\item \textbf{A fitted parameter is not automatically heritability.}
That interpretation requires valid identifying assumptions. Heritability
describes variation in a specified population and environment, not the
fraction of an individual's outcome ``caused by genes.''

\item \textbf{The Jencks critique: ``genetic'' does not mean an immutable
physiological mechanism.}
Genetic effects can operate through behavior, environments, and social
responses---for example, food preferences affecting diet and BMI. Those
pathways may be modifiable.

\item \textbf{Heritability does not determine policy effectiveness.}
A variance decomposition does not tell us what an intervention would
accomplish. Criticize Clark's specific unsupported policy conclusions while
acknowledging his stated concessions.

\item \textbf{Intergenerational mobility is not a sufficient measure of
policy success.}
Policies can improve living standards or reduce inequality without
changing relative ranks or intergenerational correlations.

\item \textbf{Genetic explanation does not establish meritocracy.}
Neither fairness nor equality of opportunity follows from evidence of
genetic influence; institutions may respond unfairly to inherited
characteristics.
\end{enumerate}


Gregory Clark's \emph{To Have and Have Not} \citep{Clark2026} makes an old
and provocative claim on the strength of an enormous new dataset: that a
person's station in life is set mainly by the genes they inherit, not by
family environment, schooling, or social policy, and that this has been true
in England, essentially without change, since 1600. The evidence comes from
two purpose-built databases: Families of England, which traces 436,000
individuals across four centuries of rare-surname lineages out to fourth
cousins, and Marriages of England, 1.7 million marriage certificates spanning
1837 to 2022.

Part I is the empirical core. Correcting for measurement error in
conventional status proxies, Clark estimates an underlying parent-child
status correlation near $0.8$ --- strong enough that its signature is still
detectable among relatives as distant as fourth cousins --- and shows this
persistence essentially unchanged across four centuries that included the
Industrial Revolution, mass education, women's enfranchisement, and the
welfare state. Across roughly a dozen kinship types, the pattern is fit by
\citet{Fisher1918}'s century-old model of additive genetic inheritance
combined with strong assortative mating, with striking apparent precision
(reported $R^2$ above $0.97$); Denmark and Sweden, despite their reputation
as high-mobility welfare states, show the same model and the same low
mobility. Clark is careful to say this is not fatalism: on his account, once social
barriers to opportunity are removed, genetic inheritance and chance become
the \emph{only} remaining sources of variation in status --- so a highly
heritable status distribution is not, by itself, evidence of an unfair
system, but potentially the signature of a meritocratic one.

Part II subjects that reading to a battery of sharp, falsifiable tests, on
the logic that direct genetic transmission --- as opposed to genes acting
indirectly through the environment parents provide --- makes specific
predictions a purely environmental account would not: mothers' and fathers'
status should matter equally (it mostly does); family size and birth order
should not matter (they mostly don't); losing a parent young, or never
meeting a grandparent, should not change a child's eventual status (it
doesn't); accidental wealth gains should dissipate within a few generations,
and regression to the mean should run at a constant rate with no wealth
traps or poverty traps (both hold). Clark counts each result as a point for
direct genetic transmission over its sociological rivals.

Part III draws out the consequences if all of this is right. Societies are,
on this view, overspending on education and welfare interventions aimed at
mobility, since mobility is set by genetic endowment and chance rather than
by policy; the only lever left for changing a population's outcomes is
changing who is in it, an implication Clark connects, without quite
endorsing it, to selective immigration; and the differential fertility of
high- and low-status parents will shape the next generation's genetic --- and
therefore social --- composition regardless of what policy does.

Whatever one concludes about that argument, the underlying data are a
genuine, durable contribution independent of it, and will outlive any
particular model fit to them. The book's insistence on the severity of
measurement error in conventional mobility proxies is also well taken and,
if anything, underappreciated outside this literature: Clark shows two
recent, careful studies of the same underlying process (long-run American
occupational mobility) reaching starkly different pictures of American
history purely as an artifact of measurement choices --- \citet{SongEtAl2020}
find high and rising mobility from 1830 to 1980, while \citet{Ward2023},
correcting for measurement error in historical occupational data, finds
persistently low mobility throughout (Fig.~4.1) --- which is exactly the
kind of problem his own attenuation-correction method is designed to
address. And the Part II battery is a well-designed set of sharp,
independently checkable predictions, whatever one makes of how any single
test comes out. None of this is in question below.

What is in question is the book's central causal claim: that this
correlational pattern specifically identifies additive genetic transmission
of status, as opposed to some other data-generating process that would fit
comparably well. That claim does not survive scrutiny, for reasons developed
in the sections that follow. Some of what follows is new; some was raised
in peer review on the underlying 2023 PNAS paper and never resolved, and the
book inherits it unchanged.

% =====================================================================
\section{The Fisher model with phenotypic assortment}
\label{sec:fisher-model}

A close fit to correlations between relatives is informative only to the
extent that we understand what generates the fitted formulas. The same
pattern of resemblance can support different explanations when different
transmission mechanisms are allowed. Spelling out the model therefore
serves two purposes: it identifies which conclusions follow from the
assumptions, and it makes clear which explanations have been excluded
before the data are examined. In particular, the absence of environmental
transmission must not be mistaken for an empirical finding that
environmental transmission is unimportant. We derive the additive,
dominance-free version of \citet{Fisher1918}'s model, following the organization of
\citet{Goldberger1978} and \citet{OttoEtAl1994}: individual components,
mating, and transmission. Measurement error is omitted throughout this
section.

\subsection{Individuals, parameters, and environmental restrictions}
\label{sec:model-components}

For person $i$, define the two column vectors
\[
 X_i=\underbrace{\begin{pmatrix}A_i\\E_i\end{pmatrix}}_{\text{random components}},
 \qquad
 q=\underbrace{\begin{pmatrix}h\\e\end{pmatrix}}_{\text{common parameters}},
 \qquad P_i=q^{\mathsf T}X_i=hA_i+eE_i.
\]
$A_i$ is the person's standardized additive genetic value; $E_i$ is the
standardized environmental component. $X_i$ varies across people, whereas
$q$ contains coefficients common to the population. The superscript
$\mathsf T$ turns a column into a row. These vectors are bookkeeping
devices, not extra components of the model. There are no dominance or other nonadditive genetic components, and no gene--environment interaction term. Assume
\begin{equation}
 \begin{gathered}
 \mathbb E[A_i]=\mathbb E[E_i]=0,\qquad
 \operatorname{Var}(A_i)=\operatorname{Var}(E_i)=1,\\
 \operatorname{Cov}(A_i,E_i)=0,\qquad
 0<h<1,\qquad e=\sqrt{1-h^2}.
 \end{gathered}
 \label{eq:model-raw}
\end{equation}
Thus $\operatorname{Cov}(X_i,X_i)=I_2$, the two-dimensional identity
matrix, and
\[
 \operatorname{Var}(P_i)=h^2+e^2+2he\operatorname{Cov}(A_i,E_i)=1.
\]
The contribution of genes on the phenotype scale is $hA_i$, with variance
$h^2$. Consequently $h^2$ is narrow-sense heritability in this population;
it is not $\operatorname{Var}(A_i)$, which we normalized to one.
$E_i$ is a substantive environmental component, not measurement error.

The sexes have the same marginal distributions, and sampling people as
parents does not change those distributions. We use the same parameters
and component variances across generations. These assumptions exclude
changes in the composition of the sampled population; stationarity of
the variances is checked below.

The model also excludes a common environmental component among biological
relatives. For the pedigrees derived below, distinct biological relatives
have $\operatorname{Cov}(E_i,E_j)=0$; an offspring's environment has zero
covariance with the components of its parents and with the genes and
environments of its siblings. We implement these exclusions formally in
the birth process below. They are additional family assumptions, not
consequences of phenotypic assortment.

\citet[pp.~10--11]{Goldberger1978} emphasizes the substantive severity of
this restriction: the classical model allows the premarital environments
of spouses to resemble one another through assortment, while excluding
environmental resemblance between siblings or between parents and
children. Shared rearing supplies no additional source of similarity;
being raised together rather than apart has no effect on the predicted
correlation for otherwise comparable biological relatives. Genes are the
only component transmitted in the model. This is an assumption about the
sources of familial resemblance, not evidence establishing their relative
importance. It does \emph{not}, however, set
$\operatorname{Cov}(E_{\rm parent},A_{\rm child})$ to zero: assortment
generates that covariance, as we show below.

\subsection{What phenotypic homogamy assumes}
\label{sec:model-ph}

Let $M$ denote the wife/mother and $F$ the husband/father. Define
\[
 \rho=\operatorname{Corr}(P_M,P_F),\qquad
 \rho_A=\operatorname{Corr}(A_M,A_F),\qquad 0<\rho<1.
\]
These definitions alone say nothing about the matching mechanism.
Our formal phenotype-only matching assumption is
\begin{equation}
 X_M\perp X_F\mid P_F,\qquad X_F\perp X_M\mid P_M.
 \tag{PH}\label{eq:model-ph}
\end{equation}
Once the wife's phenotype is known, learning her particular genetic and
environmental components gives no further information about the
distribution of the husband's components; the analogous statement holds
with husband and wife exchanged. The two statements are separate
restrictions. Symmetry of conditional independence exchanges the two
vectors but does not change the conditioning phenotype.

For clarity, $Z\perp W\mid V$ means that, conditional on $V$, the joint
probability of any event concerning $Z$ and any event concerning $W$
is the product of their conditional probabilities. Equivalently, learning
$W$ does not change the conditional distribution of $Z$ once $V$ is known.
All conditional statements hold wherever the conditional laws are
defined, up to events of probability zero. Independence is preserved by
taking functions: an event about a function of $Z$ is itself an event
about $Z$, so the same factorization applies.

We additionally assume the \emph{own-person} conditional mean is linear:
\begin{equation}
 \mathbb E[A_i\mid P_i]=hP_i.
 \tag{L}\label{eq:model-linearity}
\end{equation}
Here the prediction is of a person's genetic value from that same person's
phenotype, not from the spouse's phenotype. To see why the coefficient is
$h$, write a linear conditional mean as $\alpha+\beta P_i$. Zero means give
$\alpha=0$, and
\[
 \beta=\frac{\operatorname{Cov}(A_i,P_i)}{\operatorname{Var}(P_i)}
 =\frac{h\operatorname{Var}(A_i)+e\operatorname{Cov}(A_i,E_i)}{1}=h.
\]
A population linear projection always has this coefficient, but its
conditional mean need not equal that projection. Thus (L) is a real
assumption. From $eE_i=P_i-hA_i$ it follows that
\[
 \mathbb E[E_i\mid P_i]
 =\frac{P_i-h\mathbb E[A_i\mid P_i]}{e}
 =\frac{1-h^2}{e}P_i=eP_i,\qquad
 \mathbb E[X_i\mid P_i]=qP_i.
\]
Otto et al.'s Gaussian specification supplies such linearity. Joint
normality is sufficient here, but the spouse derivation needs only (L)
together with (PH) and the stated normalization.

\subsection{Deriving the restriction on spousal genetic resemblance}
\label{sec:model-spouses}

We spell out the conditioning steps because they separate a probability
identity from the assumptions that give the model its content.

For this argument, the four components $(A_M,E_M,A_F,E_F)$ are assumed
to have an ordinary joint density.
Because $e>0$, replacing each $E_i$ by $P_i=hA_i+eE_i$ is an invertible
change of variables, so $(A_M,A_F,P_M,P_F)$ also has an ordinary density.
We do not assume an ordinary joint density for $(A_i,E_i,P_i)$, which
contains a redundant variable.

Subscripts distinguish different density functions; lower-case letters
denote values of the corresponding random variables. Let
\[
 J=f_{A_M,A_F,P_M,P_F}(a_m,a_f,p_m,p_f),\quad
 D=f_{A_F,P_M,P_F}(a_f,p_m,p_f),\quad
 H=f_{P_M,P_F}(p_m,p_f).
\]
By the definition of a conditional density, at points with positive
denominators,
\begin{align}
 f_{A_M,A_F\mid P_M,P_F}
 &=\frac{J}{H}
   =\frac{J}{D}\frac{D}{H}\notag\\
 &=f_{A_M\mid A_F,P_M,P_F}\,
   f_{A_F\mid P_M,P_F}.
 \label{eq:model-multiplication}
\end{align}
Arguments on the last line are respectively
$(a_m\mid a_f,p_m,p_f)$ and $(a_f\mid p_m,p_f)$.
This multiplication rule is an identity; no independence has been used.

Now invoke (PH). Its second restriction, by symmetry and preservation
under functions, gives $A_M\perp(A_F,P_F)\mid P_M$, because
$A_M$ is an element of $X_M$ and $(A_F,P_F)$ is a function of $X_F$.
Conditional independence therefore factors the numerator in the
following definition:
\begin{align}
 f_{A_M\mid A_F,P_M,P_F}
 &=\frac{f_{A_M,A_F,P_F\mid P_M}}{f_{A_F,P_F\mid P_M}}\notag\\
 &=\frac{f_{A_M\mid P_M}\,f_{A_F,P_F\mid P_M}}
         {f_{A_F,P_F\mid P_M}}
 =f_{A_M\mid P_M}.
 \label{eq:model-wife-reduction}
\end{align}
The cancellation is valid on the conditional support. Similarly, the
first PH restriction gives $A_F\perp P_M\mid P_F$, and hence
\begin{align}
 f_{A_F\mid P_M,P_F}
 &=\frac{f_{A_F,P_M\mid P_F}}{f_{P_M\mid P_F}}\notag\\
 &=\frac{f_{A_F\mid P_F}\,f_{P_M\mid P_F}}{f_{P_M\mid P_F}}
 =f_{A_F\mid P_F}.
 \label{eq:model-husband-reduction}
\end{align}
Substitute \eqref{eq:model-wife-reduction} into the first factor of
\eqref{eq:model-multiplication}, and
\eqref{eq:model-husband-reduction} into the second:
\begin{equation}
 f_{A_M,A_F\mid P_M,P_F}(a_m,a_f\mid p_m,p_f)
 =f_{A_M\mid P_M}(a_m\mid p_m)
  f_{A_F\mid P_F}(a_f\mid p_f).
 \label{eq:model-density-factor}
\end{equation}
Thus, given both phenotypes, the genetic values are independent, and
each retains its own-phenotype conditional distribution. Independence
given both phenotypes \emph{alone} would not establish the latter fact.

The conditional product mean is consequently
\begin{align*}
 \mathbb E[A_MA_F\mid p_m,p_f]
 &=\int\!\!\int a_ma_f
 f_{A_M\mid P_M}(a_m\mid p_m)
 f_{A_F\mid P_F}(a_f\mid p_f)\,da_m\,da_f\\
 &=\left\{\int a_m f_{A_M\mid P_M}(a_m\mid p_m)\,da_m\right\}
   \left\{\int a_f f_{A_F\mid P_F}(a_f\mid p_f)\,da_f\right\}\\
 &=\mathbb E[A_M\mid p_m]\mathbb E[A_F\mid p_f]
   =(hp_m)(hp_f).
\end{align*}
The factorization uses PH; the final equality uses (L).
Finite second moments justify these product expectations. Finally,
average over the \emph{joint spousal phenotype distribution}, using the
law of iterated expectations:
\begin{align}
 \rho_A
 &=\operatorname{Cov}(A_M,A_F)
   =\mathbb E[A_MA_F]\notag\\
 &=\mathbb E_{(P_M,P_F)}
       \{\mathbb E[A_MA_F\mid P_M,P_F]\}\notag\\
 &=h^2\mathbb E_{(P_M,P_F)}[P_MP_F]
   =\boxed{h^2\rho}.
 \label{eq:m-def}
\end{align}
The first line uses zero means and unit genetic variances; the last uses
zero means and unit phenotype variances. The outer average is not over
independently drawn spouses.

The same functional-independence argument applies to either component
of each vector. Their conditional means are $qP_M$ and $qP_F$, so
\begin{equation}
 C:=\mathbb E[X_MX_F^{\mathsf T}]
 =\mathbb E_{(P_M,P_F)}[(qP_M)(qP_F)^{\mathsf T}]
 =\rho qq^{\mathsf T}
 =\rho\begin{pmatrix}h^2&he\\he&e^2\end{pmatrix}.
 \label{eq:ge-cross}
\end{equation}
Thus $\operatorname{Cov}(E_M,E_F)=\rho e^2$ and
$\operatorname{Cov}(A_M,E_F)=\operatorname{Cov}(E_M,A_F)=\rho he$.
The full component covariance matrix is
$\left(\begin{smallmatrix}I_2&C\\C^{\mathsf T}&I_2\end{smallmatrix}\right)$.
Here the off-diagonal blocks are matrices, not scalars; $C$ happens to
be symmetric. Since $q^{\mathsf T}q=1$, $C$ has eigenvalue $\rho$ along $q$ and zero in the perpendicular direction. The full matrix therefore has eigenvalues
$1+\rho,1-\rho,1,1$, all positive in the stated parameter range.
Also $q^{\mathsf T}Cq=\rho(q^{\mathsf T}q)^2=\rho$, verifying the
specified spouse phenotype correlation. This is Otto et al.'s
equation~(22) after deleting the cultural component.

The intuition is that spouses match on a phenotype reflecting both
components. A high phenotype predicts a genetic value of $hP$, not
$P$ itself. Applying that prediction on both sides of the marriage
multiplies phenotype covariance by $h^2$. The relationship
$\rho_A=h^2\rho$ is therefore a \emph{derived restriction}, not a third
freely chosen parameter. Without (L), PH instead gives
$\operatorname{Cov}\{\mu(P_M),\mu(P_F)\}$, where
$\mu(p)=\mathbb E[A\mid P=p]$; it need not equal $h^2\rho$.

Although this restriction appears in the equilibrium model, its
derivation requires no intergenerational argument. Stationarity of
genetic variance is a separate condition, to which we now turn.

\subsection{Births and stationarity of genetic variance}
\label{sec:model-transmission}

For child $j$ of parents $M,F$, assume additive transmission:
\begin{equation}
 A_j=\tfrac12(A_M+A_F)+S_j,\qquad P_j=hA_j+eE_j.
 \label{eq:model-transmission}
\end{equation}
$S_j$ is the segregation deviation: the child's departure from the
midparental genetic value. The coefficient $1/2$ describes additive
Mendelian transmission; it does not specify the variance of $S_j$.

Let $\mathcal H$ collect the already realized ancestral and parental
components in the pedigree before these births. For distinct births,
assume
\[
 \mathbb E[S_j\mid\mathcal H]=0,\qquad
 \mathbb E[S_jS_k\mid\mathcal H]=0\quad(j\ne k).
\]
Draw offspring environmental values independently of one another,
of $\mathcal H$, and of the segregation deviations for these births,
with mean zero and variance one. These are sufficient
no-shared-environment and no-environmental-transmission assumptions.
They concern fresh offspring environments; they do not assert
independence from mates subsequently selected on offspring phenotypes.

For any mean-zero variable $Z$ with finite variance determined by $\mathcal H$,
\[
 \operatorname{Cov}(Z,S_j)
 =\mathbb E_{\mathcal H}\{Z\,\mathbb E[S_j\mid\mathcal H]\}=0,\qquad
 \operatorname{Cov}(S_j,S_k)
 =\mathbb E_{\mathcal H}\{\mathbb E[S_jS_k\mid\mathcal H]\}=0.
\]
Here the inner means use the conditional segregation distribution given
the existing family; the outer mean uses that family's joint
distribution. These equations explain why segregation cross-terms
vanish in what follows.

Write $W=(A_M+A_F)/2$. For an arbitrary spouse genetic correlation
$\rho_A$,
\begin{align*}
 \operatorname{Var}(A_j)
 &=\operatorname{Var}(W)+2\operatorname{Cov}(W,S_j)
   +\operatorname{Var}(S_j)\\
 &=\tfrac14\{\operatorname{Var}(A_M)+\operatorname{Var}(A_F)
              +2\operatorname{Cov}(A_M,A_F)\}
   +0+\operatorname{Var}(S_j)\\
 &=\tfrac12(1+\rho_A)+\operatorname{Var}(S_j).
\end{align*}
Stationarity, $\operatorname{Var}(A_j)=1$, therefore requires
\begin{equation}
 \boxed{\operatorname{Var}(S_j)=\tfrac12(1-\rho_A)
          =\tfrac12(1-h^2\rho).}
 \label{eq:model-segregation}
\end{equation}
Random mating is the special case $\rho=\rho_A=0$, giving variance
$1/2$; zero correlation by itself need not imply independent mating
outside a Gaussian specification. Fresh $E_j$ then implies
$\operatorname{Cov}(A_j,E_j)=0$ and $\operatorname{Var}(P_j)=1$.
This is a stationary moment specification, not a proof of convergence
of a multilocus population to equilibrium. Nor does it derive $h^2$
from allele frequencies or a random-mating genetic variance.

\subsection{Parent--child and full-sibling correlations}
\label{sec:model-nuclear}

All phenotype variances equal one, so phenotype covariances below
are correlations. Start with the mother and child. From the
within-person assumptions and the spouse block $C$,
\begin{align*}
 \operatorname{Cov}(P_M,A_M)&=h,\\
 \operatorname{Cov}(P_M,A_F)
 &=h\operatorname{Cov}(A_M,A_F)+e\operatorname{Cov}(E_M,A_F)\\
 &=h(\rho h^2)+e(\rho he)
   =\rho h(h^2+e^2)=\rho h.
\end{align*}
The second line is where the PH-derived spouse covariances enter.
Using transmission and the zero segregation covariance,
\[
 \operatorname{Cov}(P_M,A_j)
 =\tfrac12 h+\tfrac12\rho h+0=\tfrac12h(1+\rho).
\]
The child's fresh environment has zero covariance with $P_M$. Therefore
\begin{equation}
 \rho_{PO}:=\operatorname{Corr}(P_M,P_j)
 =h\operatorname{Cov}(P_M,A_j)+e\operatorname{Cov}(P_M,E_j)
 =\boxed{\tfrac12h^2(1+\rho).}
 \label{eq:po}
\end{equation}
Sex symmetry gives the same formula for the father.

To expose the environmental qualification explicitly,
\begin{align*}
 \operatorname{Cov}(A_M,A_j)
 &=\tfrac12\{1+\rho h^2\}+0,\\
 \operatorname{Cov}(E_M,A_j)
 &=\tfrac12\{\operatorname{Cov}(E_M,A_M)
             +\operatorname{Cov}(E_M,A_F)\}+0
   =\tfrac12\rho he.
\end{align*}
Consequently
\[
 \operatorname{Cov}(X_M,X_j)=
 \begin{pmatrix}(1+\rho h^2)/2&0\\\rho he/2&0\end{pmatrix}.
\]
The parent's environmental component is associated with the other
parent's genes through assortment, and those genes enter the child.
No parental environment is transmitted. Goldberger makes this
distinction explicitly in his Appendix, p.~A-5. It would be incorrect
to obtain parent--child resemblance by retaining only
$h^2\operatorname{Cov}(A_M,A_j)$.

For full siblings 1 and 2, $A_1=W+S_1$ and $A_2=W+S_2$, so
\begin{align*}
 \operatorname{Cov}(A_1,A_2)
 &=\operatorname{Var}(W)+\operatorname{Cov}(W,S_2)
   +\operatorname{Cov}(S_1,W)+\operatorname{Cov}(S_1,S_2)\\
 &=\operatorname{Var}(W)+0+0+0
   =\tfrac12(1+\rho h^2).
\end{align*}
Fresh environments give
$\operatorname{Cov}(A_1,E_2)=\operatorname{Cov}(E_1,A_2)
=\operatorname{Cov}(E_1,E_2)=0$. Expanding both phenotypes,
\begin{align}
 \rho_S:=\operatorname{Corr}(P_1,P_2)
 &=h^2\operatorname{Cov}(A_1,A_2)
   +he\operatorname{Cov}(A_1,E_2)\notag\\
 &\quad+eh\operatorname{Cov}(E_1,A_2)
   +e^2\operatorname{Cov}(E_1,E_2)\notag\\
 &=\boxed{\tfrac12h^2(1+\rho h^2).}
 \label{eq:sib}
\end{align}
These are the zero-cultural-transmission, zero-shared-environment
specializations of Otto et al.'s equations~(46) and~(55).
Adding sibling environmental covariance $c$ would add $e^2c$ if the
genetic--environmental cross-covariances remained zero. The classical
formula excludes that contribution by setting $c=0$.

\subsection{Extending the pedigree}
\label{sec:model-extended}

A single mating pair does not specify how a new spouse relates to the
rest of a pedigree. To derive extended-family moments we make that
extension explicit. For each existing person $i$, draw an outside mate
$\widetilde i$ so that the pair satisfies PH and the same marginal
assumptions. Given the existing person's phenotype, the mate's
components are independent of the rest of the already realized family.
Draw distinct outside mates independently conditional on that family.
Also assume the mate-phenotype conditional mean is linear:
\begin{equation}
 \mathbb E[P_{\widetilde i}\mid P_i]=\rho P_i.
 \tag{ML}\label{eq:model-mate-linearity}
\end{equation}
The coefficient is $\rho$ because both phenotype variances are one.
Neither PH nor (L) by itself implies (ML). These additional conditions
exclude other routes connecting an incoming spouse to the pedigree.
They supply a sufficient family sampling model for the following
formulas; they were not needed for the nuclear-family results.

Let $\mathcal F$ denote the existing family before these outside mates
and their children are drawn. PH, (L), and (ML) give
\begin{align*}
 \mathbb E[A_{\widetilde i}\mid\mathcal F]
 &=\mathbb E[A_{\widetilde i}\mid P_i]\\
 &=\mathbb E_{P_{\widetilde i}\mid P_i}
       \{\mathbb E[A_{\widetilde i}\mid P_{\widetilde i},P_i]\}\\
 &=\mathbb E_{P_{\widetilde i}\mid P_i}[hP_{\widetilde i}]
   =h\rho P_i.
\end{align*}
The first equality uses the family extension, the third uses the
PH conditioning reduction and (L), and the last uses (ML).
The displayed outer average is over the mate's phenotype conditional
on the existing person's phenotype. Hence for their child $c(i)$,
\begin{align}
 \mathbb E[A_{c(i)}\mid\mathcal F]
 &=\tfrac12\{A_i+h\rho P_i\}+0\notag\\
 &=\underbrace{\tfrac12(1+\rho h^2)}_{k}A_i
   +\underbrace{\tfrac12\rho he}_{d}E_i
   =kA_i+dE_i.
 \label{eq:model-child-mean}
\end{align}
The zero term follows by averaging the zero segregation mean first
over births and then over the incoming mate. For any centered
existing-family variable $Z$, the law of iterated expectations now gives
\[
 \operatorname{Cov}(Z,A_{c(i)})
 =\mathbb E_{\mathcal F}[Z(kA_i+dE_i)]
 =k\operatorname{Cov}(Z,A_i)+d\operatorname{Cov}(Z,E_i).
\]

For a grandparent $g$, their child $i$, and grandchild $c(i)$, the
parent--child calculation gives
$\operatorname{Cov}(P_g,A_i)=h(1+\rho)/2$, whereas the fresh environment
of $i$ gives $\operatorname{Cov}(P_g,E_i)=0$. Therefore
\begin{align*}
 \operatorname{Corr}(P_g,P_{c(i)})
 &=h\{k\operatorname{Cov}(P_g,A_i)
       +d\operatorname{Cov}(P_g,E_i)\}
       +e\operatorname{Cov}(P_g,E_{c(i)})\\
 &=h\{kh(1+\rho)/2+d\cdot0\}+0
   =\boxed{\tfrac12h^2(1+\rho)k}.
\end{align*}

For an aunt/uncle $u$ and their full sibling $i$,
$\operatorname{Cov}(A_u,A_i)=k$ and all cross-component environmental
covariances between the siblings are zero. Thus
$\operatorname{Cov}(P_u,A_i)=hk$ and
$\operatorname{Cov}(P_u,E_i)=0$. Their niece/nephew satisfies
\[
 \operatorname{Corr}(P_u,P_{c(i)})
 =h\{k(hk)+d\cdot0\}+e\cdot0=\boxed{h^2k^2}.
\]

For first cousins $c(i)$ and $c(u)$, independent outside-mate draws
and the birth assumptions make their genetic values conditionally
uncorrelated given $\mathcal F$. The law of total covariance
(the sum of the mean conditional covariance and the covariance of
the conditional means) therefore gives
\begin{align*}
 \operatorname{Cov}(A_{c(i)},A_{c(u)})
 &=\mathbb E_{\mathcal F}
       [\operatorname{Cov}(A_{c(i)},A_{c(u)}\mid\mathcal F)]\\
 &\quad+\operatorname{Cov}_{\mathcal F}(kA_i+dE_i,kA_u+dE_u)\\
 &=0+k^2\operatorname{Cov}(A_i,A_u)
   +kd\operatorname{Cov}(A_i,E_u)\\
 &\quad+dk\operatorname{Cov}(E_i,A_u)
   +d^2\operatorname{Cov}(E_i,E_u)\\
 &=k^2k+0+0+0=k^3.
\end{align*}
Both cousins have fresh environments, so their phenotype correlation is
$\boxed{h^2k^3}$. This calculation exposes exactly where the additional
assumption about different incoming mates is used.

\subsection{Model restrictions, estimation, and interpretation}
\label{sec:model-tests}

Table~\ref{tab:notation} collects the moments derived above. They
match the individual-relative rows of Clark's Technical Appendix,
Table~A1, with measurement error suppressed. In Clark's notation
$r=\rho$ and $m=\rho_A=h^2\rho$; his attenuation parameter is set to
one here. Goldberger's notation has $A=\rho_A$, $c_1=h^2$, and $c_2=1$
when dominance is omitted.
\begin{table}[htbp]
\centering
\caption{Predicted phenotype correlations; $k=(1+h^2\rho)/2$}
\label{tab:notation}
\begin{tabular}{ll}
\hline
Relationship & Correlation\\
\hline
Spouses & $\rho$\\
Parent--child & $h^2(1+\rho)/2$\\
Full siblings & $h^2k$\\
Grandparent--grandchild & $h^2(1+\rho)k/2$\\
Aunt/uncle--niece/nephew & $h^2k^2$\\
First cousins & $h^2k^3$\\
\hline
\end{tabular}
\end{table}

One directly testable implication is
\begin{equation}
 \rho_{PO}-\rho_S
 =\tfrac12h^2\{(1+\rho)-(1+h^2\rho)\}
 =\boxed{\tfrac12h^2\rho(1-h^2)>0}.
 \label{eq:gap}
\end{equation}
The strict inequality uses both positive assortment and
$0<h^2<1$. Under random mating, parent--child and full-sibling
correlations both reduce to $h^2/2$; grandparent and avuncular
correlations reduce to $h^2/4$, and first-cousin correlation to $h^2/8$.
Comparisons with estimated correlations must account for sampling
uncertainty and comparable measurement and sampling across kinships.

For estimation, the no-error model has two free structural parameters,
$\theta=(h^2,\rho)$, not three independent parameters
$(h^2,\rho,\rho_A)$. For population correlations $r_\ell$ of the
relationships in the table, write
\[
 r_\ell=f_\ell(h^2,\rho).
\]
A joint minimum-distance estimator, for example, minimizes
\[
 [\widehat{\boldsymbol r}-\boldsymbol f(h^2,\rho)]^{\mathsf T}
 W[\widehat{\boldsymbol r}-\boldsymbol f(h^2,\rho)]
\]
over admissible parameters, imposing $\rho_A=h^2\rho$ by substitution.
$\widehat{\boldsymbol r}$ is the vector of estimated kinship correlations;
$W$ is a positive-definite weighting matrix. A covariance-based choice
should account for dependence between correlation estimates from
overlapping families. This formulation constrains the entire collection
of moments jointly rather than choosing each relationship separately.

\citet[p.~19]{Goldberger1978} identifies precisely the error of treating
the spousal genetic correlation as free while retaining formulas
derived under Fisher's restriction. He also notes that alternative
mating specifications may remove that restriction but then require
rederiving the kinship formulas. A broader assortment trait therefore
cannot simply be invoked to free $\rho_A$ while leaving the formulary
unexamined.

A statistically established failure of $\rho_A=h^2\rho$ rejects the
joint assumptions underlying that restriction. It does not identify
which assumption failed. Small discrepancies among estimates can be
sampling error, and a formal comparison requires a coherent unrestricted
model and its sampling covariance. Under misspecification, fitted
parameters need not identify heritability or the claimed assortment
parameters; good fit supplies no guarantee against bias. Rejection
alone does not prove that every estimate is biased, or determine the
bias's direction or magnitude. Conversely, imposing the restriction is
necessary to estimate this particular model coherently, but does not
make its environmental exclusions empirically true.

Clark's presentation obscures this distinction. On p.~16 he questions
the implication that genetic resemblance must be weaker than phenotype
resemblance for educational attainment, appealing to matching on a wider
set of characteristics. On p.~75 he instead invokes an imperfectly
measured underlying phenotype and notes that the three parameters are
linked. The latter acknowledgment does not resolve the former departure
from the model. The issue is not whether the relationship appears
somewhere in the book, but whether it is consistently treated as a
restriction required by the assumptions generating the fitted equations.

Two different arguments need to be separated. If years of education
is merely a noisy measure of the phenotype on which spouses assort,
then $\rho_A=h^2\rho$ still applies to that underlying phenotype:
$\rho$ must refer to its spousal correlation, and a measurement model
must connect its moments to those observed. If spouses instead match
on several characteristics in a way that violates PH for the modeled
phenotype, the matching model has changed. That possibility is legitimate,
but it requires specifying the new joint distribution and deriving the
resulting kinship moments. The book's appeal to broader matching does
not supply that specification or demonstrate that the same kinship
formulas remain valid. The restriction cannot be treated as an optional
extra while retaining the interpretation of formulas derived using it.

This leaves a gap between the fit and the substantive conclusion.
Clark interprets the pattern as evidence for additive genetic
transmission combined with very strong assortment on underlying social
abilities \citep[pp.~16,75,286 and Technical Appendix]{Clark2026}.
But a close fit of the correlation formulas does not, by itself,
demonstrate that this genetic mechanism---or assortment on a particular
social or cognitive trait---generated the data. That interpretation
requires a coherent specification connecting the matching process,
the transmitted components, and the observed outcomes, with its
restrictions imposed and tested. Moving between the classical model
and informal alternatives leaves the reader without a clear account
of which model the reported fit actually supports.

\paragraph{Verification.}
The accompanying verification script checks the spouse and family
covariance matrices algebraically and simulates an explicit Gaussian
construction satisfying the assumptions. It checks parent--child,
full-sibling, grandparent, avuncular, and first-cousin moments, as well
as component covariances and stationary variances. All 46 symbolic checks passed. Across nine combinations of $h^2$ and
$\rho$, simulations of 2.25 million independent pedigrees checked 1,224
component covariances and 54 phenotype moments. The largest absolute
standardized Monte Carlo discrepancies were 3.402 and 2.091, respectively,
within the prespecified six-standard-error diagnostic threshold.
The script is \texttt{verify\_fisher\_model.py}, in the repository's
\texttt{verification/} directory. The construction
provides an existence and calculation check for this moment model;
simulation is not a proof of empirical adequacy or of multilocus
equilibrium.


\section{Adding classical measurement error}
\label{sec:measurement-error}

The previous section concerns a phenotype measured without error. We now
separate that phenotype from the recorded outcome. This is a specific
measurement model: adding a residual and calling it noise is not enough
to justify multiplying every kinship correlation by the same factor.

\subsection{The observed outcome and the meaning of the attenuation factor}

Retain the standardized genetic and environmental components $A_i,E_i$,
but let the recorded outcome be
\begin{equation}
 Y_i=aA_i+eE_i+uU_i,\qquad
 X_i^+=\begin{pmatrix}A_i\\E_i\\U_i\end{pmatrix},\qquad
 w=\begin{pmatrix}a\\e\\u\end{pmatrix},\qquad Y_i=w^{\mathsf T}X_i^+.
 \label{eq:me-outcome}
\end{equation}
$U_i$ represents pure measurement error, with mean zero and variance
one. Assume $a,e>0$, $u\geq0$, and, for every pair of people $i,j$,
\begin{equation}
 \operatorname{Cov}(U_i,A_j)=\operatorname{Cov}(U_i,E_j)=0,\qquad
 \operatorname{Cov}(U_i,U_j)=0\quad(i\ne j).
 \label{eq:me-orthogonality}
\end{equation}
Thus error is uncorrelated with both true components within the person
and across people, and with other people's errors. Its variance within
the person is one, not zero. The original within-person condition
$\operatorname{Cov}(A_i,E_i)=0$ is retained; the nonzero cross-person
genetic--environmental covariances derived earlier are also retained.
These orthogonality conditions suffice for the moment calculations.
Independent errors, independent of all true family components, are a
stronger sufficient specification used in the simulations.

Standardize each recorded outcome to unit variance. The three
within-person components are uncorrelated, so
\[
 \operatorname{Var}(Y_i)=a^2+e^2+u^2=1.
\]
Define the attenuation factor and noise share as
\begin{equation}
 \eta:=u^2,\qquad
 \theta:=1-\eta=a^2+e^2>0.
 \label{eq:me-shares}
\end{equation}
Here $\eta$ is the fraction of observed variance attributable to
measurement error, whereas $\theta$ is the fraction attributable to
the true phenotype. This distinction matters for translating Clark's
notation: his $\theta$ is the \emph{attenuation factor}
\citep[p.~42 and Table~A1]{Clark2026}, so
\[
 \boxed{\theta_{\rm Clark}=\theta=1-\eta.}
\]
Thus $1-\theta=\eta$ is the noise share.

The standardized latent phenotype is
\begin{equation}
 P_i=\frac{aA_i+eE_i}{\sqrt{\theta}}
     =hA_i+e_P E_i,\qquad
 h=\frac{a}{\sqrt{\theta}},\quad
 e_P=\frac{e}{\sqrt{\theta}}=\sqrt{1-h^2}.
 \label{eq:me-latent}
\end{equation}
Consequently
\[
 Y_i=\sqrt{\theta}P_i+\sqrt{\eta}U_i,\qquad
 h^2=\frac{a^2}{\theta},\quad
 a^2=\theta h^2,\quad e^2=\theta(1-h^2).
\]
$h^2$ is heritability of the latent phenotype; $a^2$ is the additive
genetic share of variance of the recorded outcome. They coincide only
when $\eta=0$. The symbol $e_P$ denotes the environmental loading in
the error-free model; $e$ denotes its loading on the observed scale.

Assume PH and own-person conditional linearity for this same latent
$P_i$, together with the earlier transmission and pedigree assumptions.
Spouses match on $P$, not on the recording error $U$. Therefore
\begin{equation}
 \rho_A=h^2\rho=\frac{a^2}{\theta}\rho,\qquad
 \rho:=\operatorname{Corr}(P_M,P_F).
 \label{eq:me-spouse-restriction}
\end{equation}
No new transmission term is introduced: $A_j=(A_M+A_F)/2+S_j$ and
$\operatorname{Var}(S_j)=(1-h^2\rho)/2$ as before. $U$ changes what is
recorded, not what is inherited or how mates are selected.

\subsection{Why the same attenuation factor appears for individual relatives}
\label{sec:me-attenuation}

For distinct people, expanding the two observed outcomes gives
\begin{align}
 \operatorname{Cov}(Y_i,Y_j)
 &=\theta\operatorname{Cov}(P_i,P_j)
   +\sqrt{\theta\eta}\operatorname{Cov}(P_i,U_j)\notag\\
 &\quad+\sqrt{\theta\eta}\operatorname{Cov}(U_i,P_j)
   +\eta\operatorname{Cov}(U_i,U_j).
 \label{eq:me-expansion}
\end{align}
For example,
$\operatorname{Cov}(P_i,U_j)
=h\operatorname{Cov}(A_i,U_j)+e_P\operatorname{Cov}(E_i,U_j)=0$
by \eqref{eq:me-orthogonality}; the other cross-term vanishes in the
same way, and the error--error covariance is zero. Since both observed
and latent individual phenotypes have unit variance,
\begin{equation}
 \boxed{\operatorname{Corr}(Y_i,Y_j)
   =\theta\operatorname{Corr}(P_i,P_j)
   =(1-\eta)\operatorname{Corr}(P_i,P_j),\quad i\ne j.}
 \label{eq:me-attenuation}
\end{equation}
This is the source of the common multiplier in the individual-relative
rows of Clark's Table~A1. On the unstandardized signal-plus-error scale,
measurement error adds variance without adding covariance; the
correlation is therefore reduced. Attenuation occurs at the two observed
endpoints, not anew at each intervening generation.

In particular, the observed spouse correlation is
$s_Y:=\operatorname{Corr}(Y_M,Y_F)=\theta\rho$. Thus the latent
restriction may also be written
\[
 \rho_A=\frac{a^2s_Y}{\theta^2};
\]
it is not generally $a^2s_Y$. Substituting observed quantities into the
error-free restriction without the scaling factors is incorrect.
The complete spouse block for the standardized components is
\[
 \operatorname{Cov}(X_M^+,X_F^+)
 =\rho\begin{pmatrix}
 h^2&he_P&0\\he_P&e_P^2&0\\0&0&0
 \end{pmatrix}.
\]
Multiplying on the left by $w^{\mathsf T}$ and on the right by $w$
gives $\rho(ah+ee_P)^2
=\rho\{\sqrt{\theta}(h^2+e_P^2)\}^2=\theta\rho$, confirming the result.

\subsection{The adjusted moment table}
\label{sec:me-table}

Apply \eqref{eq:me-attenuation} to each moment already derived. For
example,
\[
 r_{PO}^{Y}=\theta\frac{h^2(1+\rho)}{2},\qquad
 r_S^{Y}=\theta h^2k,\qquad
 k:=\frac{1+h^2\rho}{2}.
\]
Here $r^Y$ denotes a correlation of recorded outcomes. The analogous
substitution gives the grandparent, avuncular, and first-cousin rows
in Table~\ref{tab:measurement-moments}.

The more distant collateral rows follow from the same pedigree
extension, rather than from adding further error factors.
For two collateral relatives $i,j$ whose cross-component environmental
covariances vanish, equation~\eqref{eq:model-child-mean} implies
\[
 \operatorname{Cov}(A_{c(i)},A_j)
 =k\operatorname{Cov}(A_i,A_j)+d\operatorname{Cov}(E_i,A_j)
 =k\operatorname{Cov}(A_i,A_j).
\]
Fresh offspring environments preserve the zero environmental
cross-covariances. Descending one side therefore multiplies genetic
covariance by $k$; descending both sides multiplies it by $k^2$,
using conditionally independent outside mates and births as in the
first-cousin derivation. Starting with sibling covariance $k$ gives
$k^3,k^5,k^7,k^9$ for first through fourth cousins. Removing a cousin
relationship by $t$ generations on one side supplies another $k^t$.
The observed phenotype correlation remains $\theta h^2$ times that
genetic covariance.

\begin{table}[htbp]
\centering
\caption{Correlations with classical measurement error:
$\theta=1-\eta$, $k=(1+h^2\rho)/2$}
\label{tab:measurement-moments}
\begin{tabular}{lccc}
\hline
Relationship & $n$ & $D$ & Observed correlation $r^Y$\\
\hline
Spouses & --- & --- & $\theta\rho$\\
Parent--child & 1 & 1 & $\theta h^2(1+\rho)/2$\\
Full siblings & 1 & 0 & $\theta h^2k$\\
Grandparent--grandchild & 2 & 1 & $\theta h^2(1+\rho)k/2$\\
Aunt/uncle--niece/nephew & 2 & 0 & $\theta h^2k^2$\\
First cousins & 3 & 0 & $\theta h^2k^3$\\
Second cousins & 5 & 0 & $\theta h^2k^5$\\
Third cousins & 7 & 0 & $\theta h^2k^7$\\
Fourth cousins & 9 & 0 & $\theta h^2k^9$\\
Degree-$c$ cousins, $t$ removed & $2c+1+t$ & 0
 & $\theta h^2k^{2c+1+t}$\\
\hline
\end{tabular}
\end{table}
In the final row $c\geq1$ and $t\geq0$. The columns $n,D$ specify the
regression representation below; the spouse row is a separate moment.
The single-parent and single-grandparent rows, and the collateral rows,
agree with Clark's Table~A1 after using $\theta$ as the common attenuation factor,
$r=\rho$, and $m=h^2\rho$. Averages of relatives require a different
normalization and are treated separately below.

\subsection{What the appendix regression can identify}
\label{sec:me-identification}

For the nonspousal rows, both the collateral and lineal formulas can
be written
\begin{equation}
 r^Y_{n,D}
 =\theta h^2 k^n
  \left(\frac{1+\rho}{1+h^2\rho}\right)^D,\qquad D\in\{0,1\}.
 \label{eq:me-regression-moment}
\end{equation}
For example, $n=1,D=1$ cancels $1+h^2\rho$ from $k$, leaving the
parent--child formula. $n=2,D=1$ leaves the extra factor $k$ for a
grandparent. For positive population correlations, taking logarithms gives
\begin{align}
 \log r^Y_{n,D}&=\alpha+\beta n+\gamma D,\notag\\
 \alpha&=\log(\theta h^2),\quad
 \beta=\log\left(\frac{1+\rho_A}{2}\right),\quad
 \gamma=\log\left(\frac{1+\rho}{1+\rho_A}\right).
 \label{eq:me-log-regression}
\end{align}
This is the structure of Clark's equations~A1--A3; Greek letters here
avoid confusing the regression intercept with the loading $a$.

With sufficient variation to identify all three regression coefficients,
invert them in sequence:
\begin{align}
 k&=\exp(\beta),&
 \rho_A&=2\exp(\beta)-1,\notag\\
 1+\rho&=(1+\rho_A)\exp(\gamma),&
 \rho&=2\exp(\beta+\gamma)-1,\notag\\
 h^2&=\frac{\rho_A}{\rho},&
 \theta&=\frac{\exp(\alpha)}{h^2},&
 \boxed{\eta=1-\frac{\exp(\alpha)}{h^2}.}
 \label{eq:me-inversion}
\end{align}
The crucial third line invokes the \emph{derived PH restriction}
$\rho_A=h^2\rho$. Without it, the intercept determines only
$\theta h^2=a^2$, not the separate noise and heritability shares.
Thus estimating noise by this route is possible in principle, but
its interpretation depends on the same substantive model being tested.
The exponentiated intercept is the observed genetic variance share $a^2$;
calling it latent heritability $h^2$ would be incorrect when noise is present.

The separation works because genuine $E$ enters the phenotype on which
spouses match, whereas $U$ does not. Their effects on the pattern of
parent--child and sibling resemblance therefore differ under the model.
Without that assumed distinction, two independent, untransmitted sources
of variation need not be distinguishable from these data. The recovered
noise share is consequently not an independent validation of the mating
assumptions used to identify it.

Two limits are important. With \emph{collateral relatives alone}, $D=0$:
the regression identifies $\theta h^2$ and $\rho_A$, but not $\rho$,
$h^2$, and $\eta$ separately. Even after imposing PH, different pairs
$(\theta,h^2)$ can have the same product, with
$\rho=\rho_A/h^2$, provided all parameters remain admissible.
Also, under random mating $\rho=\rho_A=0$, $\gamma=0$ for all $h^2$;
the ratio $\rho_A/\rho$ cannot identify heritability. The preceding
inversion therefore presumes positive assortment and an identified
lineal--collateral contrast. Identification can be weak near random
mating.

Estimation must enforce $0<h^2<1$, $0<\rho<1$, and
$0<\theta\leq1$. Arbitrary fitted log-regression coefficients need
not imply admissible shares. One can instead fit the three structural
parameters $(h^2,\rho,\eta)$ jointly to the moment table, imposing
the constraints directly. The spouse moment $\theta\rho$, if available,
provides an additional check on the predictions recovered from the
nonspousal regression.

Equation~\eqref{eq:me-log-regression} is exact for population moments;
using logarithms of estimated correlations introduces sampling issues.
Logs require positive estimates, and
$\operatorname{SE}(\log\widehat r)\simeq
\operatorname{SE}(\widehat r)/r$ by the delta method. For inverse-variance weighting of the log regression, the relevant
uncertainty is on the log scale. Inference must also account for
dependence between estimates from overlapping families. A small
residual sum of squares does not by itself verify the measurement or
transmission assumptions that identify $\eta$.

\subsection{Averaging parents changes the correlation}
\label{sec:me-averages}

The common multiplier $\theta$ above applies to pairs of \emph{individuals}.
It cannot be carried over mechanically to the correlation of a child
with an average of relatives. For the observed midparent value
$\overline Y=(Y_M+Y_F)/2$,
\begin{align*}
 \operatorname{Var}(\overline Y)
 &=\tfrac14\{1+1+2\operatorname{Cov}(Y_M,Y_F)\}
   =\tfrac12(1+\theta\rho),\\
 \operatorname{Cov}(Y_j,\overline Y)
 &=\tfrac12\{\operatorname{Cov}(Y_j,Y_M)
             +\operatorname{Cov}(Y_j,Y_F)\}\\
 &=\theta h^2(1+\rho)/2.
\end{align*}
Since $\operatorname{Var}(Y_j)=1$, the correlation is
\begin{equation}
 \boxed{\operatorname{Corr}(Y_j,\overline Y)
 =\frac{\theta h^2(1+\rho)}
        {\sqrt{2(1+\theta\rho)}}.}
 \label{eq:me-midparent}
\end{equation}
Even without measurement error this is
$h^2\sqrt{(1+\rho)/2}$, not $h^2$. The familiar $h^2$ is instead the
error-free \emph{regression slope} on the midparent phenotype.
With error, that slope becomes
\[
 \frac{\operatorname{Cov}(Y_j,\overline Y)}
      {\operatorname{Var}(\overline Y)}
 =\frac{\theta h^2(1+\rho)}{1+\theta\rho},
\]
which is generally not $\theta h^2$ either.

Clark's Table~A1 labels the average-parent entry a correlation and gives
$\theta h^2$. Interpreted literally as the correlation with the arithmetic
mean of the parents' outcomes, that entry does not follow from this model.
It needs a different definition or a correction. The average-grandparent
entry requires the same care: its denominator depends on the joint
covariance matrix of the four grandparents. Averaging individual-relative
correlations is not the same as correlating with an average.

\subsection{Substantive content of the common-noise correction}
\label{sec:me-meaning}

The adjustment assumes more than imperfect measurement. First, each
person's error must be uncorrelated with \emph{every} relevant relative's
true components and error. A shared informant, shared record linkage
error, or common error in assigning status to an occupation can violate
that requirement. For example, if the two error--signal terms in
\eqref{eq:me-expansion} remain zero but
$\operatorname{Cov}(U_i,U_j)=c_{ij}$, then
\[
 r^Y_{ij}=\theta r^P_{ij}+\eta c_{ij}.
\]
This is not a uniform attenuation of the original formula.

Second, the common multiplier presumes the same signal share for the
people entering each correlation. If instead
$Y_i=\sqrt{\theta_i}P_i+\sqrt{1-\theta_i}U_i$ with the analogous
orthogonality conditions, expansion gives
\[
 r^Y_{ij}=\sqrt{\theta_i\theta_j}\,r^P_{ij}.
\]
Different reliability by sex, age, cohort, outcome definition, or kinship
sample can therefore prevent cancellation in ratios. Those restrictions
must hold in the actual selected samples, not merely in a hypothetical
unselected population.

Third, the latent phenotype must be the same trait whose components
drive transmission and assortment. If an observed characteristic
affects partner choice, the part of that characteristic labeled error
relative to an unobserved trait cannot automatically be treated as
irrelevant noise. Nor does an imperfect proxy for intelligence
automatically satisfy classical measurement-error assumptions for
intelligence. A substantive difference between traits may contain
systematic genetic or environmental variation.

Finally, pure measurement error is distinct from the true environmental
component $E$. Independent, untransmitted environmental variation can
also reduce familial resemblance, and its separation from recording
error here depends on the modeled assortment and transmission
restrictions. Correcting attenuation does not remove the exclusion
of shared upbringing, validate PH, or establish a genetic explanation
for a close fit. Under a common positive $\theta$ the parent--child
minus sibling gap retains its sign:
\[
 r^Y_{PO}-r^Y_S
 =\theta h^2\rho(1-h^2)/2>0.
\]
Allowing arbitrary relationship-specific errors would change that test,
but it would also change the model being fitted.

\paragraph{Verification.}
The measurement-error verification script checks the covariance
expansion, scaling, regression inversion, and midparent normalization
symbolically. It then adds fresh measurement errors to generated
pedigrees and checks the predicted individual-relative and midparent
moments numerically. All 71 new symbolic checks passed, along with the 46 baseline checks. Across 13 parameter settings, 2.6 million simulated pedigrees passed 156 individual-relative and 26 midparent checks, including distant cousins and relatives once removed. It also checks counterexamples with unequal
reliability and correlated measurement errors. These calculations
verify the stated model; they do not verify its empirical assumptions.


\input{sections/assortment_comparison}

% =====================================================================
\section{The Identification Problem}

\subsection{Good Fit Does Not Imply Unbiased Estimates}

Set aside, for a moment, whether Clark's fit claims are even correct on their
own terms. There is a prior, more basic problem: good fit would not settle
the matter even if it held up perfectly. A correctly specified model should
generally be expected to fit the data well and to yield unbiased parameter
estimates. The converse does not hold. A misspecified model can fit
observed correlations closely while its parameter estimates mean nothing
close to what the author claims they mean, because the fitted parameters are
absorbing the specification error along with the genuine signal. Treating
close agreement between predicted and observed correlations as validation of
the assumptions used to derive the predictions --- rather than as one
necessary, but far from sufficient, condition for those assumptions to be
right --- is a fallacy with a long and specifically relevant history in
exactly this empirical literature.

\subsection{The 1970s Precedent}
\label{sec:1970s-precedent}

In the mid-1970s two rival groups --- Jinks and Eaves at Birmingham, and Rao,
Morton, and Yee at Honolulu --- fit competing biometric models of essentially
the same kind to overlapping British and American IQ kinship-correlation
data, and each group reported an excellent fit. Goldberger's (1978)
re-examination of both efforts, transcribed in full in the companion file
\texttt{Goldberger\_1978\_ModelsAndMethods.tex} included in this repository,
is instructive on two counts.

First, the two models, despite each fitting its own data well, produced
substantively incompatible conclusions. Fit to overlapping American data,
Rao and Morton's own estimate of narrow heritability of IQ among children was
$.69$, against Jinks and Eaves's $.34$ --- roughly a twofold disagreement on
the same underlying quantity, from ostensibly similar exercises. The two
groups also disagreed about the sign of the adult-versus-childhood pattern:
Rao and Morton's estimate of broad heritability among American adults was
$.30$, falling relative to children, while Jinks and Eaves's corresponding
estimate was $.68$, unchanged from their childhood estimate. A reader told
only that both models ``fit the data well'' would have no way to know that
the two exercises disagreed this sharply about what the data implied.

Second, and more pointed, Goldberger showed algebraically that the Honolulu
group's original 8-parameter model was not even \emph{identified}: ``the
model was indeterminate, that is `non-estimable,' that is `not identified':
even if the population values of all 11 kinship correlations were known,
unique values of the [model's] parameters could not be obtained.'' Five of
the eight parameters could take on any of an entire family of values without
changing a single one of the model's predictions. Rao, Morton, and Yee
themselves eventually conceded the indeterminacy and published a
``corrected'' version of the model that forced identification by simply
fixing one of the indeterminate parameters at an arbitrarily chosen value.
Goldberger then computed that this patched model's predicted correlations
still matched the data closely (``the fit is clearly very good,''
$\chi^2_4=2.96$). The sequence is the point: a model can be shown
mathematically not to pin down its own parameters, be patched by fixing one
of them by fiat, and still return an excellent-looking fit --- fit is simply
not informative about whether the underlying model was identified in the
first place, let alone correctly specified.

\subsection{The Same Pattern, on Clark's Own Data}

The restriction derived in Section~\ref{sec:model-tests} provides a concrete
way to examine Clark's own numbers. Table~\ref{tab:table45} reproduces
the raw sibling and parent-child (``Full Sibling'' and ``Child'') correlations
from the book's Table 4.5, for the two outcomes with two separate birth-year
windows: occupational status and higher-education attainment.

\begin{table}[h]
\centering
\caption{Raw correlations from Clark's Table 4.5}
\label{tab:table45}
\begin{tabular}{lcccc}
\hline
 & \multicolumn{2}{c}{Occupational status} & \multicolumn{2}{c}{Higher education} \\
 & 1780--1859 & 1860--1919 & 1780--1859 & 1860--1919 \\
\hline
Full Sibling ($\hat\rho_S$)      & 0.584 & 0.496 & 0.558 & 0.359 \\
Child ($\hat\rho_{PO}$)          & 0.585 & 0.513 & 0.500 & 0.305 \\
\hline
$\hat\rho_{PO}-\hat\rho_S$       & $+0.001$ & $+0.017$ & $-0.058$ & $-0.054$ \\
\hline
\end{tabular}
\end{table}

If a common, positive, phenotype-specific attenuation factor scales
both correlations for a given outcome and period, it preserves the sign
of the parent--child minus sibling difference. Equation~\eqref{eq:gap}
requires that population difference to be positive when assortment is
positive and $0<h^2<1$. The higher-education point estimates have the
opposite sign in both periods. This is a useful specification diagnostic,
but point estimates alone do not establish rejection: a formal test must
account for their joint sampling uncertainty and verify comparable
measurement and sampling across relationship types. If attenuation differs
between the two kinships, the simple sign comparison need not be valid.
The relevant exercise is therefore a joint test of the constrained
predictions, not an inference from overall fit alone.

\subsection{External Evidence: Population versus Direct Genetic Effects}

A further way to adjudicate between observationally similar models --- short
of the kind of kinship-correlation moment-matching this section has just
shown to be unreliable on its own --- is to bring in evidence external to the
correlations being fit. The clearest example concerns the value of $m$ (the
genotypic spousal correlation) that Clark treats as consistent with modern
genetic evidence. Genome-wide association study estimates of the kind Clark
cites for parental genetic similarity are estimates of \emph{population}
effects, not \emph{direct} (causal) effects, and a substantial recent
literature shows these need not be proportional to one another
\citep{YoungEtAl2019,YoungEtAl2022}. Treating an $m\approx0.6$ derived from
population-genetic estimates as evidence about \emph{direct} genetic
transmission of status --- the quantity the book's causal claims actually
rest on --- is not obviously licensed once that distinction is taken
seriously. Related work on ``genetic nurture,'' in which a parent's genotype
affects a child's outcomes through the environment the parent provides rather
than through the genes transmitted to the child \citep{KongEtAl2018}, is
exactly the kind of indirect environmental-transmission channel that the
no-environmental-transmission assumption of Section~\ref{sec:model-components} rules out by fiat.

The book's own Chapter 6 takes up genetic nurture directly and concludes that
recent attempts to demonstrate its importance fail, arguing that twin and
adoption studies would already have detected such effects if they were large.
This does not engage the population-versus-direct-effects distinction above,
which is a separate concern from the twin/adoption evidence Chapter 6
addresses; nor does it engage \citet{KongEtAl2018}'s design, which detects
genetic nurture using a method (non-transmitted parental alleles) that does
not rely on the twin/adoption comparisons Chapter 6 critiques. I raised a
version of this gap directly in my referee report on the underlying PNAS
paper, where Clark's response offered an unpublished personal simulation
arguing that direct-effects-only models with sufficient assortment could
reproduce apparent genetic-nurture patterns; that simulation was not, and has
not since been, made available for independent scrutiny.

\subsection{Confounding Within Monozygotic Twin Pairs}

A related, longer-standing literature makes the same kind of point about
within-family comparisons more generally. MZ-twin-difference designs, widely
used to estimate the return to schooling on the logic that identical twins
share both genotype and family background \citep{Griliches1979}, are only as
good as the assumption that within-pair schooling differences are unrelated
to individual ability \citep{Griliches1977,BoundSolon1999,Neumark1999}.
Directly testing this assumption, \citet{SandewallEtAl2014} find that
within-pair differences in adolescent IQ predict within-pair differences in
educational attainment among monozygotic twins, and that controlling for IQ
reduces the within-pair estimated return to schooling by roughly 15\%. If
ability differences of exactly this kind survive even within monozygotic
twin pairs, treating any specific empirical design --- Clark's collateral-
relative decay curve included --- as if it cleanly isolates a single
additive genetic channel, free of confounding by within-family differences
in unmeasured ability, is not something that can simply be assumed.

% =====================================================================
\section{What ``Genetic'' Does Not Mean}

\subsection{Genes Fixed Is Not the Same as Genes' Effects Fixed}

Even granting Clark's estimates at face value, the inferential leap from
``status is highly heritable'' to ``society has little room to change social
outcomes'' does not go through. ``Genetic,'' in the variance-decomposition
sense used throughout this model, means only that an outcome is statistically
associated with genotype; it says nothing about the causal pathway by which
that association arises, and that pathway can run through channels ordinary
usage would call environmental and modifiable. Body-mass index is highly
heritable by exactly this kind of estimate --- classic twin studies find
substantial resemblance even among monozygotic twins reared apart in
different homes \citep{StunkardEtAl1990} --- yet the pathway from genotype to
body mass runs substantially through food choice and health behavior --- both
manifestly alterable. The point, due to Christopher Jencks, is that a
genotype being fixed (nobody chooses their own genes) is a completely
different claim from a genotype's \emph{effects} being fixed or immutable.
Reared-together versus reared-apart comparisons already discussed in
Section~\ref{sec:model-components} make the same point from the environmental side: if shared
environment truly contributed nothing, as the model's sole-genetic-channel
assumption requires for every relationship except marriage, those comparisons
should show no difference, and they do.

\subsection{Heritable Is Not the Same as Unchangeable}

The book goes further than the general inference above and states the
stronger claim directly: ``The second implication of direct genetic
transmission is that social outcomes will change \ldots\ only when the
genetic endowment of the society changes'' (p.~10). This is not a defensible
reading of what heritability estimates license, and standard examples from
outside the social sciences make the point sharply. Phenylketonuria is,
absent intervention, about as close to a purely genetic determinant of
outcome as exists in human biology --- yet it is entirely preventable by a
change in diet, an intervention that has nothing to do with the genetic
endowment of anyone. Myopia is highly heritable and is fully correctable with
eyeglasses. Closer to the social outcomes this book studies, a literacy
program that teaches a population to read can change welfare enormously
regardless of the heritability of whatever cognitive or social trait
``literacy'' loads onto --- heritability describes the sources of variation
\emph{within} a given environmental range, not a ceiling on what a new
environmental intervention outside that range can do.

\subsection{Ordinal Mobility Is Not the Only Thing Policy Can Be For}

There is also a narrower, more specific problem with how the book gets from
its estimates to its policy conclusions. The Introduction's own overview of Part III states:
``The first
implication of the importance of direct genetic transmission of social
outcomes is that modern societies spend too much on education \ldots\ What
drives much of this expenditure is a belief that boosting years of education
by less advantaged students increases social mobility rates. Yet despite a
substantial narrowing of the gap in years of education between high and low
status students 1851--2022, social mobility rates are unchanged'' (p.~9). The
argument silently assumes that the \emph{only} legitimate justification for
education spending is changing relative, ordinal mobility --- moving people
up and down the status distribution relative to one another, which is what
the book's intergenerational transmission coefficients measure. But a policy
can be worth its cost by raising everyone's \emph{absolute} outcomes even if
it leaves the ordinal ranking, and the estimated transmission coefficients,
completely untouched: a literacy program that lifts the reading level of an
entire cohort improves welfare regardless of whether the same families end up
in the same relative positions afterward. Evidence that an intervention did
not change the intergenerational correlation is evidence about mobility; it
is not, by itself, evidence about whether the intervention was worthwhile.

The specific empirical claim that follows --- ``studies of the effect of
increases in compulsory years of education for less able students in
1918--22, 1947, and 1972 find no effect on employment, income, or health''
(p.~9) --- is also not uncontested. Using the same 1972 reform (the raising of
the minimum school-leaving age in England, Scotland, and Wales from 15 to
16), \citet{BarcellosEtAl2018} find that the additional year of compulsory
schooling it induced had real effects on middle-age health: it reduced body
size and improved lung function (with larger improvements on both margins
for individuals with higher genetic predisposition to obesity), and also
increased blood pressure. Whatever one makes of the mixed direction of these
effects, they are effects --- directly on the health margin the book says
these reforms show ``no effect'' on.

Nor is the birth-order test in Part II as distinctively confirmatory of the
genetic model as the book presents it: prior labor-economics literature
already finds birth-order effects on educational and economic outcomes to be
real but small \citep{BlackDevereuxSalvanes2005}, so a modest or null
birth-order effect in the FOE data is also compatible with ordinary
environmental accounts of status transmission, not only with the additive
genetic model.

% =====================================================================
\section{Conclusion}

A pattern recurs across the specific problems catalogued above: evidence that
would seem to cut against the genetic-transmission story is consistently
absorbed as further confirmation of it, rather than treated as a live threat
to the model. Strong marital assortment in the English data is read as
confirming the model. When a far richer Swedish extended-kinship dataset,
spanning 141 distinct relationship types, is used by \citet{ColladoEtAl2023}
to reject the additive genetic model outright on its own terms, the
rejection is attributed not to a real specification failure but to Collado,
Ortu\~no-Ort\'in, and Stuhler's failure to correct for measurement error in
years of education --- even though, by that same logic, essentially all of
Clark's own point estimates would need to be treated with equal suspicion,
not only the ones that happen to be inconvenient. And the near-equality of
the English parent-child and sibling correlations is invoked in the book as
an argument against a rival social-transmission model, although it also
warrants checking against the Fisher model's parent--child versus sibling
restriction (Section~\ref{sec:model-tests}), with sampling uncertainty and
measurement assumptions taken into account. In each case, the same body of
evidence would be read one way if it favored the model and the opposite way
if it did not, and the book consistently finds the model favored.

None of this need reflect bad faith, and it is not a new phenomenon. It is a
documented occupational hazard of exactly this kind of model-fitting
exercise, which is why this review has tried to root the critique
historically rather than treat it as a one-off lapse specific to this book:
the same asymmetry --- treating a close fit as confirmation and treating a
poor fit as someone else's data problem --- is what let two incompatible
research programs each claim vindication from the same 1970s IQ-kinship
data, and it is what let a model later shown not even to be identified pass
as ``clearly very good'' once patched (Section~\ref{sec:1970s-precedent}). The lesson from that
history was that goodness of fit, on its own, says nothing about correct
specification. \emph{To Have and Have Not} does not appear to have absorbed
that lesson.

None of this should obscure what the book gets right, and it is worth
restating plainly rather than as a rhetorical concession. The genealogical
infrastructure --- Families of England, Marriages of England --- is a
genuine and lasting contribution to the study of long-run mobility,
independent of any conclusion drawn from it. The book's insistence that
measurement error in conventional status proxies is severe enough to distort
mobility estimates is a real methodological point, well illustrated, and
likely to outlast the debate over its author's preferred model. And the
falsifiable test battery in Part II is good scientific practice, whatever
one makes of how any individual test comes out.

What has not been established is the specific causal claim the book is
organized around: that this pattern of kinship correlations identifies
additive genetic transmission of status, as opposed to some other
data-generating process --- including some the book never considers --- that
would fit comparably well. That is a strong, specific, and separable claim,
and on the evidence presented it has not met its own evidentiary bar.

I want to be explicit about what kind of disagreement this is. I am not
troubled by the hypothesis that genes influence social outcomes, and I do
not regard raising that possibility as objectionable, dangerous, or a mark
against the author's character, in the way much of the criticism this book
will attract is likely to suggest. My disagreement is narrower, and I hope
more useful, than that: on the specific evidence and the specific model
presented here, I am not persuaded.

\bibliographystyle{apalike}
\bibliography{references}

\clearpage
\section*{Working Notes: Points to Revisit}
\begin{itemize}
\item \textbf{Parent and grandparent averages in Table A1.}
Revisit the erroneously labelled average-parent and average-grandparent
entries in the book's technical appendix. Distinguish a correlation with
an average, a regression slope on that average, and an average of
individual-relative correlations. Averaging changes reliability and the
normalization, so the common individual attenuation factor cannot simply
be carried over. Check the practical relevance separately: these average
rows do not appear among the Table 4.5 moments specified for estimation.
\item \textbf{Intuition for nonidentification of measurement error under
genetic assortment.} Explain why independent true environmental variation
and independent recording noise both dilute the observed genetic signal
without affecting mate selection on genetic value. The correlations reveal
only the product $\theta h^2$: less reliable measurement of a more
heritable trait is observationally equivalent to more reliable measurement
of a less heritable trait. Even the observed spouse correlation does not
separate the two in this model.
\item \textbf{The fundamental measurement-error exclusion.}
An imperfect measure of a latent trait is not automatically that trait
plus classical measurement error. The gap between observed education and
an asserted latent social ability may reflect family resources, access
to schooling, preferences, discrimination, and sex-specific opportunities.
These influences may covary with the person's own genetic and environmental
components, relatives' components, and relatives' measurement residuals.
Our common-$\theta$ correction sets these covariances to zero. If they
are nonzero, additional terms enter the observed kinship covariances and
a uniform attenuation factor need not apply. Develop this objection:
an admissible estimate of $\theta$ does not validate the exclusions or
establish that an error-free heritability has been recovered.
\item \textbf{Identification with generation-specific measurement error.}
Revisit identification when attenuation differs across generations rather
than being governed by one common $\theta$. How can the difference between
full-sibling and parent--child correlations then distinguish the lineal
factor $(1+r)/(1+m)$ from different attenuation factors? The same question
arises for the single-grandparent moment relative to the aunt/uncle moment.
Specify whether reliability varies by historical generation or by the
number of generations separating a pair, derive the resulting pair-level
attenuation factors, and establish what additional restrictions or moments
would separate measurement error from the effects of $r$ versus $m$.
\item \textbf{Measurement correction added to A1--A3, but parameter recovery unclear.}
The book inserts $\theta$ into the moment equations, so the A3 intercept
identifies $\exp(a)=\theta h^2$, not latent heritability alone. Yet the
surrounding text still describes estimated heritability and cites the PNAS
SI Table S2. Check whether the plotted/reported quantity is $\theta h^2$
or a separately recovered $h^2$. Under phenotypic assortment, recovery
from A3 is possible using $m=2b-1$, $r=(1+m)\exp(c)-1$, $h^2=m/r$,
and $\theta=\exp(a)/h^2$, subject to admissibility; retaining the same
regression is not itself an identification error. Distinguish this
passage from Chapter 5's use of Table A2's spouse and affine-relative
moments for Denmark and Sweden. Under the stated common-attenuation
formulas, the sibling-in-law/sibling ratio identifies $r$, while the
spouse correlation identifies $\theta r$, providing an additional route
to separate reliability and heritability. Audit that estimation and
Table 5.4 separately rather than assuming the PNAS procedure applies.
\item \textbf{Reconcile Figure A1, PNAS Figure 1, and PNAS SI Table S2.}
The book's Figure A1 (p.~298) is the counterpart of PNAS Figure~1,
but several plotted coordinates differ. Approximate graphical readings
of the heritability axis put higher education 1860--1919 near 0.58 in
PNAS and 0.33 in the book, and occupational status 1860--1919 near
0.66 and 0.53, respectively. PNAS Figure~1 also disagrees with its own
SI Table S2: occupational status 1780--1859 is plotted near 0.94,
whereas the table reports 0.722. These are graphical readings, not
recovered numerical estimates. Establish the source data, estimation
version, and parameter interpretation for each display; do not infer
from the changed points alone that measurement error or the structural
constraint was newly handled correctly. For an independent extension,
the book's Table 5.1 (p.~67) reports observed education correlations for
spouses and several affine relatives in Denmark and Sweden. Keep these
distinct from the English education cohorts and from the latent
correlations inferred in Table 5.2.
\item \textbf{An apparent impasse between assortment interpretations.}
With phenotypic assortment and free measurement attenuation, the lineal
contrast and $m=h^2r$ allow separate recovery of $h^2$ and $\theta$.
However, our unweighted A3 fits to Table 4.5 yield inadmissible structural
recoveries for four of eight outcomes: $r>1$ for company directorship,
and $h^2>1$ for occupational status 1780--1859, higher education
1780--1859, and literacy. Alternatively, if the intended assumption is
genotypic assortment, the moment formulas must change: the separate
lineal adjustment disappears ($c=0$), and these kinship correlations
identify only $m$ and $\theta h^2$. Intergenerational observations can
improve estimation and test the common decay pattern, but do not add
the identifying contrast needed to separate $h^2$ from $\theta$.
Thus one cannot rescue the interpretation merely by switching assortment
assumptions while retaining separately identified heritability. Is this
an impasse for the interpretation of the reported estimates? Distinguish
that concern from a formal rejection: bounded phenotypic estimation
remains possible, the results depend on weighting and sampling uncertainty,
and externally justified information about reliability or heritability
could resolve the genotypic model's nonidentification.
\item \textbf{Claims about regression to the mean.}
Otto et al. explicitly discuss linear regression toward the mean in
cultural models, including the possibility of reproducing the genetic
model's rate. Clark is therefore wrong to present linear regression
toward the mean as a distinct prediction of genetic transmission.
Distinguish the shape of the regression relationship from its numerical
rate, and distinguish a fixed prediction from a fitted parameter: in the
genetic benchmark at issue the coefficient is 0.5, whereas in the cultural
model it is a free parameter, which can also take that value. This is a
difference in predictive restriction, not a demonstration that cultural
transmission cannot generate the same pattern. Supply the explicit
passages from Otto et al. and Clark, and specify the assumptions behind
the genetic 0.5 benchmark rather than treating it as universal.
\item \textbf{Key points to hammer home.}
There is genuine value in Clark's analyses, and the review should
acknowledge that explicitly. The central concern is not that the analyses
are without merit, but whether the conclusions he draws follow logically
from them. Preserve this distinction when assessing both contributions
and limitations.
Often Clark's conclusions do not follow from the evidence he presents.
This does not establish that the conclusions are wrong; it means they
have not been established nearly as credibly as he suggests. Develop
this distinction consistently throughout the review. Working analogy
to consider: his apparent resistance to alternative explanations recalls,
in the opposite direction, the stubbornness of Lysenkoite social
scientists who instinctively reject evidence allowing any role for
genetics. The target is asymmetric standards of evidence, not an
equivalence of substantive views. He also sometimes exaggerates how
strongly his tests discriminate between genetic and cultural models;
symmetry and regression toward the mean are key examples to document.
Examine his claims that school reforms yield no benefits, distinguishing
evidence of no benefit from failure to establish a benefit. Track shifts
between policy criteria: does a reform help anyone or improve average
outcomes, does it help the poor, does it reduce inequality, and does it
increase or reduce mobility? These are distinct questions, and a failure
on one criterion does not establish failure on the others. Collect
specific passages showing where the criterion changes within an argument.
\item \textbf{Selective discussion of assumptions in twin and adoption studies.}
Clark's discussion focuses on one assumption whose violation can bias
estimated $h^2$ downward, but these designs depend on other assumptions
as well. Identify the assumption he emphasizes and examine the others,
including violations that could bias estimates upward or have an
ambiguous effect. Distinguish twin from adoption designs and assess the
direction of each bias under explicit conditions. A possible source of
downward bias alone does not establish that the estimates are lower
bounds on heritability.
\item \textbf{Lower resemblance when raised apart and claims about adoptees.}
Add a discussion of the evidence for lower resemblance among relatives
raised apart than among those raised together, and what this comparison
implies for Clark's interpretation of environmental influences. His
conclusions about adoptees appear much stronger than the evidence
warrants. Identify the precise claims and supporting studies, report
the relevant comparisons and uncertainty, and distinguish what those
studies establish from broader claims about the absence of effects of
rearing environments.
\item \textbf{Within-MZ-twin estimates are not automatically causal.}
Acknowledge Clark's valid point that comparisons within monozygotic
twin pairs do not by themselves identify causal effects. Controlling
for shared genes and family background does not eliminate nonshared
confounding or measurement error. Also note that he does not cite the
prior literature documenting these limitations, either within economics
or outside it. Supply the relevant references and check his discussion
and bibliography when developing this point: the methodological concern
is important, but it is not new.
\item \textbf{Implications for education spending and social outcomes.}
Discuss the claims: ``First implication\ldots\ is that modern societies
spend too much on education.'' ``What drives much of this expenditure
is a belief that boosting years of education by less advantaged students
increases social mobility rates.'' ``Second implication\ldots\ social
outcomes will change only when the genetic endowment of the society
changes.'' This implies that maximum genetics explains 60\%.
\item \textbf{The covariance exclusions matter beyond genetic nurture.}
Revisit Clark's rejection of genetic nurture on p.~86, but do not frame
the objection as depending on this particular causal mechanism. Any
relevant covariance that the fitted model incorrectly sets to zero can
bias the recovered genetic, assortment and attenuation parameters.
This includes environmental covariance between relatives and
cross-person genetic--environmental covariance beyond the terms already
implied by the stated assortment model. Establishing or rejecting a
specific ``nature of nurture'' pathway does not establish the required
zero-covariance restrictions. Write out the omitted covariance terms and
show how they enter the fitted moments; do not presume that every
violation biases every parameter, or that the direction is always the
same. The issue is the validity of the identifying exclusions, not the
label attached to the source of resemblance.

\item \textbf{Binary phenotypes and the scale of the model.}
Examine whether the linear phenotype, kinship-correlation and common
measurement-error formulas apply to binary outcomes such as higher
education, literacy and company directorship. Distinguish correlations
and heritability on the observed binary scale from those of a latent
liability. Investigate whether prevalence, thresholds, and differences
across cohorts or sexes alter the predicted moments, attenuation and
identification. Work out the implications before treating a binary
outcome as a continuous phenotype with classical measurement error.
\item \textbf{Fit competing cultural and genetic transmission models.}
Assess the absence of a serious comparison with joint cultural--genetic
transmission or a purely cultural model. A good fit of the restricted
genetic model, even if established, does not identify that mechanism.
Specify and compare the alternatives under transparent assumptions and
the same data and objective, distinguishing predictive restrictions
from extra flexibility. For the regression-to-the-mean point above,
\citet[pp.~8--9, eq.~8]{OttoEtAl1994} write
\[
 B_o=\beta_I(B_m+B_f)+\delta S_C,
\]
where $S_C$ is cultural transmission error, not schooling. The equation
resembles genetic transmission, but $\beta_I$ is not fixed at one half.
They explicitly describe regression towards the mean under cultural
transmission; this pattern alone cannot select the genetic explanation.
\item \textbf{Disconnect between behavioural genetic estimates and policy conclusions.}
Develop the distinction between variance-component estimates and the
causal effects of feasible interventions. Fixed genes do not imply
fixed effects, and a statistical direct genetic effect need not operate
through an unmodifiable physiological pathway. Use the eyeglasses
example and relate the point to \citet{Goldberger1978} and
\citet{Jencks1980}. Give explicit credit for the provocative questions,
valuable dataset and interesting patterns while explaining why the
model-selection and policy conclusions go beyond what those results
establish. Engage with the empirical evidence rather than personal labels.

\item \textbf{Assortment on ``something deeper'' requires a model.}
Address Clark's argument that people assort on something deeper than
the observed phenotype. Specify what that latent characteristic is,
how it relates to observed outcomes and their genetic and environmental
components, and how it governs partner selection. Derive the spouse
covariances and kinship moment conditions under that matching rule,
including the measurement assumptions needed to connect latent and
observed outcomes. An appeal to a deeper trait is not sufficient to
retain the original moment formulas or parameter interpretations.
If this is simply phenotypic assortment on a latent trait measured with
classical error, state and justify those restrictions explicitly; if it
is a different assortment mechanism, establish which moments and
identification results change.

\item \textbf{Selective treatment of bias and lower within-family GWAS estimates.}
Examine whether Clark acknowledges and adequately addresses the much
lower estimates reported in within-family GWAS. His emphasis on biases
that could depress other estimates should be assessed alongside evidence
that controlling for family background reduces population GWAS effect
estimates or polygenic-score associations. Identify the relevant studies,
report the comparable within-family and population estimates with their
uncertainty, and check his discussion before asserting an omission.
Distinguish SNP effects, score associations or prediction, and SNP-based
heritability from total narrow-sense heritability; these are not
interchangeable quantities. Assess explanations involving indirect
genetic effects, population structure, assortment and measurement error
under explicit assumptions, rather than treating the lower estimates
as either automatically decisive or automatically dismissible.

\item \textbf{Use within-family PGIs in spousal analyses.}
Clarify that analyses intended to interpret spousal PGI resemblance as
assortment on direct genetic effects should use PGIs constructed with
within-family GWAS effect estimates. Population-GWAS weights can also
reflect population structure and indirect genetic associations, so
spousal correlations in those scores need not isolate the intended
genetic component. Compare the results using population and within-family
weights, allowing for differences in precision and predictive accuracy.
Within-family weighting is not a complete solution: establish what the
score measures, which biases remain, and how its spousal correlation
relates to the model's latent additive-genetic correlation $m$.

\item \textbf{Misspecification need not produce poor fit.}
Make explicit that a misspecified model can fit the observed moments
well: fitted parameters may absorb omitted covariance terms, and
different mechanisms may generate similar patterns of resemblance.
Consequently, good fit does not validate the identifying restrictions
or the genetic interpretation of the estimates. Keep this logical point
separate from the empirical concern that fit may itself be poor once
the model's restrictions and admissible parameter bounds are imposed.
Assess both, but do not make the identification criticism depend on
establishing poor fit or a formal statistical rejection.

\item \textbf{Residual stratification despite ancestry restrictions and quality control.}
Document the substantial GWAS evidence that population stratification
can bias estimates even after standard quality control and restriction
to an apparently homogeneous ancestry group. Such restrictions do not
establish the absence of fine-scale population structure or its
association with social and environmental differences. Assemble concrete
examples, specifying the phenotype, ancestry definition, controls and
evidence used to diagnose residual bias. Explain the implications for
interpreting population-GWAS effects and spousal PGI correlations as
purely genetic quantities. Distinguish stratification from indirect
genetic effects and assortative mating; a within-family versus population
difference alone need not identify which mechanism accounts for it.

\item \textbf{The Collado--Clark debate and the identification disputes of the 1970s.}
Explore the parallels with the debates involving Jencks, the Birmingham
group, and Cloninger and colleagues. A recurring problem is identification
of distinct transmission mechanisms from a small set of kinship moments:
similar fits can conceal very different causal models and parameter
interpretations. Check the historical positions and the precise points
of disagreement before drawing the comparison. The irony is how much
progress has occurred in data, molecular measurement and research designs,
while this basic identification problem can reappear when inference rests
on a limited set of resemblance patterns. Explain which newer evidence
can actually distinguish the competing mechanisms, rather than treating
technical progress itself as a solution to identification.

\item \textbf{Modern evidence on dominance does not establish zero shared environment.}
Develop the point that modern literature supports Clark's position on
dominance more than he acknowledges. Identify the relevant findings,
state exactly which claim about dominance they support, and delimit
the outcomes and designs to which they apply. This does not vindicate
the separate claim that shared-environment variance $c^2$ is zero for
all outcomes in adulthood. Evaluate that claim against outcome-specific
evidence and uncertainty, keeping dominance and shared-environment
components conceptually distinct. Support for one modelling assumption
cannot stand in for evidence for the other.
% User's dictated continuation after "or" is still pending; do not infer it.
\end{itemize}

\end{document}
